%! The Empirical Rule — Unit 4, Lesson 4.5 — Lesson Plan
% EFFL plan (experience first, formalize later). Teacher-facing — the spoiler rule does NOT
% apply here: use the formal vocabulary and the lettered standard codes freely.
% This course HAS `fixedskillbox` (unlike AP Statistics): use it where a tabularx must stay
% intact on one page, and plain `skillbox` + \boxguard everywhere else.
\documentclass[10pt]{article}
\usepackage{saar-article}
\usepackage{saar-boxes}
\usepackage{graphicx}   % -article does NOT load graphicx; needed for thumbnails/screenshots
\graphicspath{{images/}}
\newcommand{\TallMath}[1]{$\displaystyle #1\rule[-1.4em]{0pt}{3.2em}$}
\newcommand{\UnitNumberName}{Unit 4: Random Variables and the Normal Distribution \quad}
\newcommand{\LessonNumberName}{Lesson 4.5: The Empirical Rule}

\begin{document}
\begin{center}
    {\Large\bfseries \CourseName \\
    {\small\bfseries \UnitNumberName \LessonNumberName}}
\end{center}

\begin{tcolorbox}[colback=frost, colframe=cerulean, boxrule=0.9pt, arc=2mm,
    left=2mm, right=2mm, top=1.5mm, bottom=1.5mm]
    \textbf{Primary Objective:} TODO --- what students can do, interpret, and \emph{justify}
    by the end.
    \\[2pt]
    % CITE THE STANDARD CODES. This is a standards-driven course; the codes appear here and
    % again on the Close & Preview / Connections line.
    \textbf{Standards:} TODO (e.g.\ PS.DS.1a, PS.DS.1c)
    \\[2pt]
    \textbf{Lesson model:} \emph{Experience first, formalize later.} Students work the activity
    in small groups using only prior knowledge; the teacher circulates with questions, cues, and
    prompts (not answers); the debrief attaches the formal vocabulary to what the groups already
    found; the application uses it once together, in context. No separate direct-instruction
    block, no tiered instruction, no exit ticket.
\end{tcolorbox}

\renewcommand{\arraystretch}{1.4}
\begin{skillbox}[Priority Ideas \& Skills]{goldbox}
% Fill in the \item text ONLY. Two tabularx cells, one itemize each — do not re-nest or
% re-balance the environments (that is the classic source of a stray \end{itemize}).
\begin{tabularx}{\linewidth}{@{}X|X@{}}
\textbf{Learning targets (student language)} & \textbf{Key understandings (the ``why'')} \\[2pt]
\hline
\vspace{-6pt}
\begin{itemize}[leftmargin=1.1em, nosep, after=\vspace{2pt}]
  \item TODO: ``I can \dots'' target, traceable to a lettered Knowledge \& Skill (cite the code).
\end{itemize}
&
\vspace{-6pt}
\begin{itemize}[leftmargin=1.1em, nosep, after=\vspace{2pt}]
  \item TODO: the why --- and state the lesson's \textbf{target misconception} explicitly.
\end{itemize}
\end{tabularx}
\end{skillbox}

\begin{skillbox}[{Vocabulary, Concepts, \& Theorems}]{greenbox}
    % TODO: key terms + definitions for this lesson. The debrief attaches these; they must
    % NOT appear on the cover, in the warm-up, or in the pre-debrief slides (spoiler rule).
\end{skillbox}

% fixedskillbox never splits, so the thumbnail and its caption stay together.
\begin{fixedskillbox}[Activate Prior Knowledge \& Spiral Review]{frost}
    \begin{tabularx}{\linewidth}{@{}X c@{}}
        \begin{minipage}[t]{0.55\linewidth}\vspace{0pt}
            % TODO: the ~3 warm-up items and the SEED each plants (which formal idea it sets up
            % WITHOUT naming it), plus what to debrief aloud. Resist naming the vocabulary here.
        \end{minipage}
        &
        \begin{minipage}[t]{0.38\linewidth}\vspace{0pt}\centering
            % TODO: spiral-review thumbnail. Authored warm-ups compile to
            % target/, so leave this text-only unless you keep a source
            % PDF in the warmup/ directory to embed.
        \end{minipage}
    \end{tabularx}
\end{fixedskillbox}

\begin{fixedskillbox}[Lesson at a Glance --- \MeetingLength]{frost}
\renewcommand{\arraystretch}{1.25}
\begin{tabularx}{\linewidth}{@{}l c X X@{}}
\textbf{Phase} & \textbf{Min} & \textbf{Students} & \textbf{Teacher} \\
\hline
Warm-Up & 5 & TODO & Circulate; debrief item~3 aloud \\
Experience \& Formalize: Activity & 22 & TODO activity title, scenarios 1--2, groups of 3--4 & Circulate; questions, cues, prompts --- \emph{not answers}; choose work to display \\
Debrief: Formalize & 14 & Share findings; record QuickNotes & Name each idea on the groups' work (in a second color); fill QuickNotes together \\
Application & 10 & TODO application title, worked together; compute, interpret, justify & Lead it; students hold the pen, you hold the questions \\
Close \& Assign & 4 & Record the assignment; preview next lesson & Name what changed today; choose packet CYU or DeltaMath \\
\end{tabularx}
\end{fixedskillbox}

\boxguard
\begin{skillbox}[Experience \& Formalize --- The Activity: \emph{TODO Title}]{frost}
\textbf{Launch (2 min).} TODO: the launch script --- the data context, the scenarios, and the
charge: \emph{``Use what you already know. For every answer, be ready to show me where you see
it in the data.''}

\begin{multicols}{2}
\textbf{What students do.} TODO: the arc of the two scenarios --- what scenario~1 builds, what
scenario~2 varies, and which item is the \textbf{crux question} (the target misconception).

\columnbreak
\textbf{What the teacher does.} Circulate. Listen before speaking. When a group stalls, give a
\emph{question, cue, or prompt}, not an answer:
\begin{itemize}[leftmargin=1.1em, nosep]
  \item TODO: item ref: ``TODO prompt.''
\end{itemize}
While circulating, pick the group work to display in the debrief.
\end{multicols}
\end{skillbox}

\boxguard
\begin{skillbox}[Debrief --- Formalize (what goes on the board, in a second color)]{frost}
Work through the displayed student work in order, and \textbf{attach each formal name to something
a group already wrote}. Students record each line in their QuickNotes as you go.
\begin{multicols}{2}
\begin{enumerate}[leftmargin=1.3em, nosep, itemsep=3pt]
  \item TODO: red-ink move --- ``student phrase'' $\to$ \emph{formal term / notation}.
\end{enumerate}
\columnbreak
\textbf{QuickNotes.} TODO: the walkthrough of the packet's QuickNotes box --- a summary of what
the groups discovered, not a new lecture.

\textbf{Why this order.} TODO --- including any example the teacher poses cold (the special case
the activity does not carry).
\end{multicols}
\end{skillbox}

\boxguard
\begin{skillbox}[Application --- \emph{TODO Title} (worked together, 10 min)]{frost}
TODO: the one problem worked together --- the first place the just-named vocabulary is \emph{used},
and the lesson's only in-class practice. It carries the whole arc: state the question in context,
compute in the packet's \texttt{work} block, \textbf{interpret the result in context}, then the
``what if we change a number?'' part and the \textbf{justification} it demands.
\textbf{Run it with the students holding the pen and you holding the questions:} TODO the three
questions you ask.

\textbf{This is the formative check.} There is no exit ticket --- read the Interpret and What-if
responses over shoulders as you circulate and sort them into TODO categories to decide how the
next lesson opens.
\end{skillbox}

\boxguard
\begin{skillbox}[Watch For (while circulating)]{redbox}
\begin{itemize}[leftmargin=1.2em, nosep]
  \item TODO: misconception, keyed to an item number.
\end{itemize}
Cold-call prompts: TODO.
\end{skillbox}

\boxguard
\begin{skillbox}[Homework --- Check Your Understanding (assigned, scored)]{goldbox}
TODO: the assignment --- $\sim$6 CYU-style items in new contexts, spanning this lesson's lettered
skills and spiraling one earlier idea, each ending in an interpretation or a justification. This
is the lesson's practice set and it \emph{is} scored (the cover's score column carries a blank
for it). TODO: how it is scored and when it is due.

\textbf{Packet or DeltaMath --- decide this lesson.} The packet set is authored either way. If
DeltaMath carries adequate content for TODO (this lesson's skills), assign the DeltaMath set
instead and tell students to keep the packet pages as their worked reference; otherwise assign
the packet pages. Say which one aloud at Close \& Assign so it is unambiguous.
\end{skillbox}

\boxguard
\begin{skillbox}[Close \& Preview]{goldbox}
\textbf{Assign the practice} (name packet or DeltaMath) and name what changed today rather than
closing on the assignment alone.
\textbf{Connections:} TODO --- the lettered standard codes this lesson covered, and where they go next.
\textbf{Preview:} TODO next lesson.
\end{skillbox}

% ── Teacher notes — one per component, in packet order ──
% TEACHERNOTES: teacher prose lives HERE, never in a _key. A note in a key is the one block
% with no counterpart in the blank, so it makes the key run longer and costs the student
% packet a blank page.
\begin{teachernote}[Warm-Up]
    % TODO: pacing + what each item seeds; what to cut if it runs long.
\end{teachernote}

\begin{teachernote}[Experience \& Formalize]
    % TODO — \textbf{Activity.} pacing split across the two scenarios; the early-finisher move.
    % \textbf{Debrief.} the must-land moments; what to cut if time is short.
    % \textbf{Application.} the questions to ask while the students hold the pen; where the
    % \texttt{work} block is; how to sort the formative responses.
\end{teachernote}

\begin{teachernote}[Homework]
    % TODO: the CYU items and what each targets; the scoring approach; whether DeltaMath
    % covers this lesson's skills; which item to go over at the start of the next lesson.
\end{teachernote}
\end{document}
